\clearpage\corrigeDevoir{6}
\begin{blocLivre}{dscor-6-1}\begin{corrigeBac}{6}{Un circuit série}
\demarche{Le courant est commun aux deux résistances ; leurs tensions s’ajoutent.}
\schema[corrige]{serie}
\begin{questions}
\item $9,0=(100+200)I$, donc $I=0,030$ A.\pts{2 pt}
\item $U_1=100\times0,030=3,0$ V ; $U_2=6,0$ V.\pts{2 pt}
\item Ampèremètre en série, voltmètre en dérivation sur $R_2$ ; couper l’alimentation avant de rebrancher.\pts{2 pt}
\end{questions}
\end{corrigeBac}\end{blocLivre}
\begin{blocLivre}{dscor-6-2}\begin{corrigeBac}{7}{Une acquisition de température}
\demarche{Passer successivement du code à la tension, à la résistance puis à la température.}
\schema[corrige]{pont-capteur}
\begin{questions}
\item $U=5,00\times500/1000=2,50$ V.\pts{2 pt}
\item $R_c=10,0\times2,50/2,50=10,0$ kΩ ; $\theta=25,0$ °C.\pts{2 pt}
\item Non : la température calculée reste inférieure à 30 °C.\pts{1 pt}
\item Vérifier le nouvel étalonnage et le domaine de température, ainsi que les tensions admissibles de la carte et les connexions.\pts{2 pt}
\end{questions}
\end{corrigeBac}\end{blocLivre}
\begin{blocLivre}{dscor-6-3}\begin{corrigeBac}{7}{Une solution au laboratoire}
\demarche{Rassembler préparation quantitative, montage et contrôle d’une transformation.}
\schema[corrige]{dilution}
\begin{questions}
\item $V_p=c_fV_f/c_m=10,0$ mL ; pipette de 10,0 mL et fiole de 100,0 mL.\pts{2 pt}
\item Le réfrigérant condense les vapeurs, qui retombent dans le ballon ; arrivée d’eau en bas, sortie en haut, montage ouvert à l’atmosphère.\pts{2 pt}
\item Non : cela soutient l’identité de l’espèce détectée, mais des impuretés peuvent ne pas être séparées ou révélées.\pts{1 pt}
\item La hausse suggère un caractère exothermique ; la température dépend aussi des pertes, des quantités et du récipient.\pts{2 pt}
\end{questions}
\end{corrigeBac}\end{blocLivre}

